% Esquisse d'un programme (cote 119, copy 1): archive cote 119, site page 51 (PDF page 52), typescript page 49.
% typescript. Draft compared: EsquisseFr.pdf (text layer):1764-1816.
% AI-transcribed by Claude Opus 5.5 (claude-opus-5-5, adaptive thinking, effort high), Message Batches API, batch esq1 (msgbatch_01LNsogFyF3y6KSp9UbtJtBk), 2026-11-12; not yet checked by eye.
\scanpage{119/51}
\mspage{49}

bien souvent, à des obstacles psychiques considérables. C'est pour-%
quoi, à certains égards et de plus en plus ces dernières années,
mes activités d'enseignant ont souvent agi comme un poids, plu-%
tôt que comme un stimulant pour le déploiement d'une réflexion
géométrique tant soit peu avancée, ou seulement délicate.

9. L'occasion me semble propice ici de faire un bref bi-%
lan de mon activité enseignante depuis 1970, c'est-à-dire depuis
que celle-ci s'effectue dans un cadre universitaire. Ce contact
avec une réalité très différente a été pour moi riche en ensei-%
gnements, d'une portée d'un tout autre ordre d'ailleurs que sim-%
plement pédagogique ou scientifique. Ce n'est pas ici le lieu de
m'étendre sur ce sujet. J'ai dit aussi au début de ce rapport le
rôle qu'a joué ce changement de milieu professionnel dans le
renouvellement de mon approche des mathématiques, et celui de mes
centres d'intérêt en mathématique. Si par contre je fais le bilan
de mon activité enseignante au niveau de la recherche proprement
dite, j'aboutis à un constat d'échec clair et net. Depuis plus
de dix ans que cette activité se poursuit an par an au sein d'une
même institution universitaire, je n'ai pas su, à aucun moment,
y susciter un lieu où ``il se passe quelque chose'' - où quelque
chose ``passe'', parmi un groupe si réduit soit-il de personnes, re-%
liées par une aventure commune. A deux reprises, il est vrai, vers
les années 74 à 76, j'ai eu le plaisir et le privilège de susciter
chez un élève un travail d'envergure, poursuivi avec élan : chez
Yves Ladegaillerie le travail signalé précédemment (par. 3) sur
les questions d'isotopie en dimension 2, et chez Carlos Contou-Carrère
(dont la passion mathématique n'avait pas attendu la rencontre
avec moi pour éclore) un travail non publié sur les jacobiennes
locales et globales sur des schémas de bases généraux (dont une
partie a été annoncée dans une note aux CR). Ces deux cas mis à
part, mon rôle s'est borné, au cours de ces dix ans, à transmettre
tant bien que mal des rudiments du métier de mathématiciens à quel-%
ques vingt élèves au niveau de la recherche, ou tout au moins à

\margin{51}
